\subsection{Ausgangslage}
\label{subsec:ausgangslage}

Die Einrichtung neuer Softwareprojekte erfordert wiederkehrende Entscheidungen
zu Architektur, Technologien, Schnittstellen, Konfiguration sowie Test-,
Build- und Deployment-Prozessen. Werden diese Grundlagen für jedes Vorhaben
unabhängig erstellt, entstehen unterschiedliche technische Ausgangslagen.
Web-, Cloud- und Desktop-Anwendungen stellen zudem verschiedene Anforderungen
an Komponenten und Bereitstellung. Bei plattformübergreifenden Projekten
müssen gemeinsame Bestandteile deshalb klar von spezifischen Erweiterungen
getrennt werden.

Eine standardisierte technische Basis kann wiederkehrende Entscheidungen
bündeln und ein konsistentes Entwicklungsmodell bereitstellen. Gegenstand der
Fallstudie ist das Projekt \enquote{Template-Projekte}, ein wiederverwendbares
Full-Stack-Technologie-Template für Web-, Cloud- und Desktop-Projekte. Seine
Ausgestaltung und Eignung werden im weiteren Verlauf systematisch untersucht
\parencite{kleiveist2026template}.
